% -----------------------------------------------------------------------
\section{Applications}
\label{sec:applications}
% -----------------------------------------------------------------------

Once trained, the surrogate is used as a fast forward model. We consider three
uses: autonomous inference (rollout), uncertainty quantification, and ensemble
data assimilation.

\subsection{Inference}
\label{sec:inference}

At inference the trained surrogate is run autonomously: starting from an initial
state $\state_0$ it is advanced by repeatedly applying the surrogate step
(Eq.~\ref{eq:step}), feeding each prediction back as the input for the next step
(Eq.~\ref{eq:rollout}), with the inflow parameters supplied at every step. Because
the geometry mask is re-applied at each step, obstacle cells remain exactly zero
throughout the rollout.

\paragraph{Ensemble initialisation via training-data warm start.}
\label{sec:warmstart}
A critical practical challenge in surrogate-based inference and data assimilation
is the
initialisation of the ensemble state. Starting members from zero velocity or
a simple mean-flow profile requires an extended spin-up period before the
surrogate produces physically consistent turbulent flow statistics, consuming
assimilation windows during which the observations are not yet informative.

We address this with a \emph{training-data warm start}: each ensemble member is
initialised from the final velocity snapshot of a randomly selected \palm{} training
trajectory. Because these snapshots are genuine LES equilibrium states, the
surrogate begins its rollout from a physically consistent condition and no
spin-up period is required. The inflow parameters for each member are
simultaneously anchored to the inflow conditions that prevailed at the end of the
selected training trajectory, preserving the physical consistency between the
initial velocity field and the inflow forcing.

Crucially, the truth observations against which the ensemble is assimilated are
generated by \palm{} (the full LES solver), while the forward model used in the
\esmda{} loop is the neural surrogate. This \emph{non-identical} pairing avoids the
\emph{inverse crime} that would arise if the same model generated both the
observations and the forecast ensemble, and provides a realistic assessment of the
surrogate's ability to track a high-fidelity truth trajectory.

\subsection{Uncertainty Quantification}
\label{sec:uq}

% TODO: describe surrogate-based uncertainty quantification (e.g. ensemble
% propagation of the inflow-parameter prior through the surrogate to obtain
% predictive velocity-field uncertainty, calibration diagnostics). Placeholder.
\textit{This section is under preparation: it will describe how predictive
uncertainty is quantified by propagating the inflow-parameter prior through the
surrogate ensemble and assessing the calibration of the resulting spread.}

\subsection{Data Assimilation}
\label{sec:esmda}

\begin{figure}[t]
  \centering
  \begin{tikzpicture}[font=\small, x=1cm, y=1cm]

%% ---------------------------------------------------------------
%% INITIALISATION COLUMN  (x = 0)
%% ---------------------------------------------------------------

\node[surblock, fill=priorblue!30, text width=1.8cm, align=center]
  (init_warm) at (0, 1.0)
  {\textbf{Warm start}\\ $\state_0^{(n)}\!\sim\!\mathcal{D}$};

\node[surblock, fill=priorblue!20, text width=1.8cm, align=center]
  (init_prior) at (0, -0.4)
  {AR(2) prior\\ $\bth_0^{(n)}\!\sim\!p(\bth)$};

\begin{scope}[on background layer]
  \node[draw=black!40, dashed, rounded corners=6pt, inner sep=7pt,
        fit=(init_warm)(init_prior)] (init_box) {};
\end{scope}
\node[font=\small\itshape] at (0, 2.95) {Initialisation};


%% ---------------------------------------------------------------
%% WINDOW 1  (x = 3.5)
%% ---------------------------------------------------------------

\node[font=\small\bfseries] at (3.5, 2.95) {Window 1};

\node[netblock, fill=netcolor!25, align=center,
      minimum width=1.8cm, text width=1.6cm, minimum height=1.1cm]
  (w1_rollout) at (3.5, 1.6)
  {\textbf{P3D}\\ Rollout};

\foreach \dy in {-0.25,-0.08,0.08,0.25}{
  \draw[stateblue!60, line width=0.5pt, -{Latex[length=2pt,width=1.5pt]}]
    ([yshift=\dy cm]w1_rollout.west) --
    ([yshift=\dy*0.4 cm]w1_rollout.east);
}

\node[star, star points=5, fill=obsyellow, draw=obsyellow!70!black,
      minimum size=0.5cm, inner sep=1pt]
  (w1_obs) at (3.5, 0.3) {};
\node[above=1pt of w1_obs, font=\tiny] {$\yobs_w$};
\node[right=3pt of w1_obs, font=\tiny, align=left]
  {6 sensors\\ $z{=}3\,$m};

\node[surblock, fill=postcolor!30, text width=1.8cm, align=center]
  (w1_esmda) at (3.5, -0.6)
  {ES-MDA\\ $\textstyle\sum\tfrac{1}{\alpha_i}{=}1$};

\node[surblock, fill=postcolor!50, align=center]
  (w1_post) at (3.5, -1.6)
  {Posterior $\xctrl^a_w$};

\draw[myarr] (w1_rollout.south) -- (w1_obs.north)
  node[midway, right=2pt, font=\tiny] {forecast};
\draw[myarr] (w1_obs.south) -- (w1_esmda.north);
\draw[myarr] (w1_esmda.south) -- (w1_post.north);


%% ---------------------------------------------------------------
%% WINDOW 2  (x = 7.0)
%% ---------------------------------------------------------------

\node[font=\small\bfseries] at (7.0, 2.95) {Window 2};

\node[netblock, fill=netcolor!25, align=center,
      minimum width=1.8cm, text width=1.6cm, minimum height=1.1cm]
  (w2_rollout) at (7.0, 1.6)
  {\textbf{P3D}\\ Rollout};

\foreach \dy in {-0.25,-0.08,0.08,0.25}{
  \draw[stateblue!60, line width=0.5pt, -{Latex[length=2pt,width=1.5pt]}]
    ([yshift=\dy cm]w2_rollout.west) --
    ([yshift=\dy*0.4 cm]w2_rollout.east);
}

\node[star, star points=5, fill=obsyellow, draw=obsyellow!70!black,
      minimum size=0.5cm, inner sep=1pt]
  (w2_obs) at (7.0, 0.3) {};
\node[above=1pt of w2_obs, font=\tiny] {$\yobs_w$};
\node[right=3pt of w2_obs, font=\tiny, align=left]
  {6 sensors\\ $z{=}3\,$m};

\node[surblock, fill=postcolor!30, text width=1.8cm, align=center]
  (w2_esmda) at (7.0, -0.6)
  {ES-MDA\\ $\textstyle\sum\tfrac{1}{\alpha_i}{=}1$};

\node[surblock, fill=postcolor!50, align=center]
  (w2_post) at (7.0, -1.6)
  {Posterior $\xctrl^a_w$};

\draw[myarr] (w2_rollout.south) -- (w2_obs.north)
  node[midway, right=2pt, font=\tiny] {forecast};
\draw[myarr] (w2_obs.south) -- (w2_esmda.north);
\draw[myarr] (w2_esmda.south) -- (w2_post.north);


%% ---------------------------------------------------------------
%% CONTINUATION  (x = 10.0)
%% ---------------------------------------------------------------

\node[font=\large] at (10.2, 0.0) {$\cdots$};


%% ---------------------------------------------------------------
%% INTER-COLUMN CONNECTING ARROWS
%% ---------------------------------------------------------------

%% Init -> Window 1 rollout  (both arrive at w1_rollout.west staggered)
\draw[myarr] (init_warm.east)  -- ++(0.25,0) |- (w1_rollout.west);
\draw[myarr] (init_prior.east) -- ++(0.25,0) |- (w1_rollout.west);

%% Window 1 posterior -> Window 2 rollout
\draw[myarr] (w1_post.east) -- ++(0.3,0) |- (w2_rollout.west);

%% Window 2 posterior -> continuation
\draw[myarr] (w2_post.east) -- ++(0.4,0) |- ++(0.6,0.4) -- ++(0.4,0);


%% ---------------------------------------------------------------
%% COLUMN SEPARATOR LINES
%% ---------------------------------------------------------------

\draw[black!20, line width=0.5pt, dashed] (1.8, 3.3) -- (1.8, -2.1);
\draw[black!20, line width=0.5pt, dashed] (5.3, 3.3) -- (5.3, -2.1);
\draw[black!20, line width=0.5pt, dashed] (8.8, 3.3) -- (8.8, -2.1);

\end{tikzpicture}

  \caption{\esmda{} workflow with the \pthreed{} neural surrogate. Ensemble members are
    initialised from \palm{} training snapshots (warm start) and rolled forward by
    the surrogate over each assimilation window. Multiple \esmda{} analysis passes
    update both the flow state and inflow parameters conditioned on point
    observations. The posterior initialises the next window.}
  \label{fig:esmda}
\end{figure}

We embed the surrogate in the \esmda{} scheme of Section~\ref{sec:esmda_framework}.
The assimilation proceeds over multiple sequential windows. At the start of
each window the surrogate propagates each ensemble member forward in time,
accumulating predictions at the observation times. \esmda{} then performs
$\niter$ analysis passes to update the ensemble, and the resulting
posterior initialises the forecast ensemble for the next window. In the
experiments of Section~\ref{sec:da_results} the augmented control vector is the
time-varying inflow-parameter trajectory $(\iang, \ivel)$, with the velocity
field propagated by the surrogate conditioned on the updated parameters.

\subsubsection{Observation Operator}
\label{sec:obs}

Observations are provided by $6$ virtual point sensors arranged on a $3 \times 2$
spatial grid at street level ($z = 3\,\mathrm{m}$), located in open street
canyons where the flow is physically accessible. Each sensor observes all
three velocity components $(u, v, w)$ via temporal mean aggregation over fixed
assimilation intervals, yielding $18$ scalar observations per assimilation window.
Observation errors are modelled as isotropic Gaussian noise, with a common
variance that reflects the combined effect of sensor noise and representativeness
error. The observation operator $\Hobs$ extracts the ensemble-member
predictions at the sensor locations and averages them over the window duration,
mirroring the aggregation applied to the synthetic truth observations.
